\documentclass[11pt]{article}
\usepackage[margin=0.9in]{geometry}
\usepackage{amsmath,amssymb,amsthm,booktabs,mathrsfs}
\usepackage[T1]{fontenc}
\usepackage{lmodern,microtype}
\usepackage[colorlinks=true,linkcolor=blue,urlcolor=blue]{hyperref}
\newtheorem{theorem}{Theorem}[section]
\newtheorem{lemma}[theorem]{Lemma}
\newtheorem{proposition}[theorem]{Proposition}
\newtheorem{corollary}[theorem]{Corollary}
\theoremstyle{remark}
\newtheorem{remark}[theorem]{Remark}
\newcommand{\norm}[1]{\left\|#1\right\|}
\title{Toward a four-dimensional Yang--Mills construction\\
\large Notebook 24: inherited interactions and the next-step norm test}
\author{Research working document}
\date{17 October 2026}
\begin{document}
\maketitle
\begin{abstract}
We separate an inherited analytic quantum correction into its
stationary-background value and its fluctuation increment in
the actual ordinary pivot chart. Directional Cauchy bounds the
increment by a factor of order $gR$ divided by the available
analytic margin, without a carrier-size derivative loss.
The exact pivot incidence confines its projected support to
the original blocks and their relevant endpoint neighbors.
Connected-carrier compression then yields an explicit
sufficient inequality for the stronger interaction norm
required by the next Gaussian integration. Under that
inequality and the existing covariance hypotheses, the
inherited contribution to the connected logarithm has the
same small factor. The inequality is not yet verified for
the actual output decay constants as the block side varies.
This result makes that parameter test, and the independent
domain-stability problem, explicit.
\end{abstract}

Notebooks 21--23 provide the classical-reference integral,
stationary-background transport, and local gauge covariance.
The following argument treats the additional interactions
produced by earlier integrations. It keeps their full
background term outside the integral and controls the
remaining fluctuation effect under a displayed norm budget.


% Conditional inherited-interaction bridge. No preamble. Prefix qi-.
\section{A conditional bridge for inherited quantum interactions}
\label{qi-section}

This section is a finite-step result with an explicit unverified
weight inequality. It uses the actual source fluctuation chart and
does not claim that the available constants meet that inequality,
that field domains are stable under iteration, or that an RG
induction has closed. Fix odd $L\ge5$, color and compatible periodic source
geometry; all displayed constants are uniform in old refinement
$M$ and volume, but may depend on $L$.

\subsection{The exact split before the next integration}

After stationary transport, the activities are functions of the
next-current unit-spacing potential, now denoted $a$. Their
connected carriers $H$ are unit current cells. The following
integration groups these cells into $L$-block Gaussian cells,
with dimensions
\begin{equation}\label{qi-dimensions}
 D_G=d(L^d-1)\dim\mathfrak g,\qquad
 D_{\mathrm{bg}}=dL^d\dim\mathfrak g.
\end{equation}
The first counts retained Gaussian coordinates; the second counts
all background link coordinates in a block. The inherited unit-cell
activity has at most $D_0=d\dim\mathfrak g$ background scalar
coordinates per unit cell.

Write $V_*(W)=e^{ia_*(w)}$, $W=e^{iw}$, for the actual matched
stationary branch. The exact source insertion is
\begin{equation}\label{qi-actual-insertion}
 U=e^{i\Phi_{V_*(W)}(g\xi)}V_*(W).
\end{equation}
Keep its current background independent initially, and define
\begin{align}
 F(a,z)&=\frac1i\log\left(e^{i\Phi_{e^{ia}}(z)}e^{ia}\right),
 &F(a,0)&=a,                                            \label{qi-map}\\
 \Delta E_H(a,\xi)&=E_H(F(a,g\xi)|_H)-E_H(a|_H).         \label{qi-increment}
\end{align}
Then, on the common small domain,
\begin{equation}\label{qi-exact-split}
 \sum_H E_H\left(\left.\frac{\log U}{i}\right|_H\right)
 =\sum_H E_H(a_*(w)|_H)
       +\sum_H\Delta E_H(a_*(w),\xi).
\end{equation}
The first sum is outside the fluctuation integral. It is not
discarded as a vacuum constant unless it is independent of $w$.
There is no factor $g^{-2}$ in \eqref{qi-increment}: $E_H$ is
already a dimensionless quantum exponent. Its linear and
quadratic fluctuation terms may remain bounded interactions.

\subsection{A value estimate without a carrier-size derivative loss}

Suppose $E_H$ is bounded holomorphic on the scalar-coordinate
product ball $|a|<b$, with $e_H=\norm{E_H}_b$, and
\begin{equation}\label{qi-inherited-norm}
 \sup_i\sum_{H\ni i}e^{a_0|H|}e_H\le\epsilon,
 \qquad a_0>0,\qquad H\text{ connected}.
\end{equation}
The audited fixed-$L$ source chart supplies
$b_{\mathrm{ch}},z_{\mathrm{ch}},C_{\mathrm{ch}}>0$ such that
\begin{equation}\label{qi-displacement}
 \norm{F(a,z)-a}_\infty\le C_{\mathrm{ch}}\norm z_\infty,
 \quad \norm a_\infty<b_{\mathrm{ch}},\quad
       \norm z_\infty<z_{\mathrm{ch}}.
\end{equation}
All fixed coordinate norm conversions are included in these constants.
Choose $b_-<\min(b,b_{\mathrm{ch}})$ and set
\begin{equation}\label{qi-domain-budget}
 m=b-b_->0,\quad |g|R<z_{\mathrm{ch}},\quad
 \delta=C_{\mathrm{ch}}|g|R<m,\quad
 u=\frac{\delta}{m-\delta}.
\end{equation}

\begin{lemma}[Uniform inherited increment]\label{qi-increment-lemma}
On the independent background ball $|a|<b_-$ and retained
window $\norm\xi_\infty\le R$,
\begin{equation}\label{qi-increment-bound}
 \norm{\Delta E_H}\le u e_H.
\end{equation}
The increment is holomorphic in the independent background.
If $\delta\le m/2$, then $u\le2C_{\mathrm{ch}}|g|R/m$.
\end{lemma}
\begin{proof}
For fixed allowed $(a,\xi)$, put
$h=F(a,g\xi)|_H-a|_H$. The one-variable function
$f(\zeta)=E_H(a|_H+\zeta h)$ is bounded by $e_H$ on
$|\zeta|<m/\delta$. Its Taylor series gives
\[
 |f(1)-f(0)|\le e_H\sum_{n\ge1}(\delta/m)^n=u e_H.
\]
The actual chart is holomorphic in $a$, proving the final statement.
\end{proof}
This estimates one directional value increment, not the absolute
sum of its multivariate Taylor coefficients. It therefore has
no $D_0|H|$ factor and does not replace the coordinate-count
argument needed for strict support substitution. No cutoff
indicator is differentiated.

\subsection{Exact pivot geometry after Gaussian-cell projection}

Use positive current bonds, assigned to their initial site, and
let $\pi_L$ be centered $L$-cube projection. A nonpivot output
of the actual map satisfies
\[
 \Phi_V(z)(b)=z(b),\qquad
 F(a,z)(b)=\frac1i\log(e^{iz(b)}e^{ia(b)}),
\]
so it reads only that bond. A pivot $p(c)$ for
$c=(y,y+Le_\mu)$ lies on the common interface of the two
endpoint blocks, with initial site in the $y$ block. Its exact
implicit equation reads only the retained links and background
links in these two blocks. No other pivot enters the equation,
so solving it does not propagate dependence farther.

Define
\begin{equation}\label{qi-projected-carrier}
 Y(H)=\pi_L(H)\cup
 \{\text{other endpoint block of }c: p(c)\text{ starts in }H\}.
\end{equation}
The increment reads only the independent $(a,\xi)$ fibers in
$Y(H)$. This carrier is connected, and, with
$\chi=d+1$ and $N_{\mathrm{near}}=2d+1$,
\begin{align}
 |Y(H)|&\le\chi|\pi_L(H)|,
 &\ell(Y(H))&=|Y(H)|-1,                              \label{qi-projected-size}\\
 \mathbf1_{\{B\in Y(H)\}}
 &\le\sum_{i\in H}\mathbf1_{\{d_{\mathrm{block}}(B,\pi_L(i))\le1\}}.
                                                               \label{qi-projected-anchor}
\end{align}
Indeed every extra block attaches to $\pi_L(H)$, and at most
$d$ outward endpoint blocks are added per projected block.
The conservative neighbor count $N_{\mathrm{near}}$ is enough
for anchoring. A convention reading both incoming and outgoing
endpoints may instead use $\chi=2d+1$. The positive oriented
source convention gives the displayed $d+1$.

This uses the actual one-pivot incidence. Enlarging each input
cell by an entire radius-$O(L)$ current ball and paying its
$O(L^d)$ volume would lose the sharper estimate.

\subsection{The explicit sufficient interaction budget}

Group equal carriers into the exact exponent
\begin{equation}\label{qi-grouped-exponent}
 Q_Y(a,\xi)=\sum_{H:Y(H)=Y}\Delta E_H(a,\xi).
\end{equation}
The absolute anchored estimates below justify this grouping.
The strict-support, tunable real-baseline theorem of Notebook18
requires the multiplier norm
\begin{align}
 \epsilon_{\mathrm{int}}
 &=\sup_B\sum_{Y\ni B}v_Y\mathcal K_\theta^{|Y|}
                       e^{4|Y|+\kappa\ell(Y)},
 &v_Y&=\sup_{\mathrm{cutoff}}|e^{Q_Y}-1|,              \label{qi-integration-norm}\\
 \mathcal K_\theta
 &=\left(\frac{1+\theta}{(1-\theta)^3}\right)^{D_G/2},
 &0<\theta&\le1/4.                                    \label{qi-density-cost}
\end{align}
The cutoff-local multiplier is
$\chi_Y(e^{Q_Y}-1)$, extended by zero to all real fluctuations.
All additional source restrictions remain present.

Put
\begin{equation}\label{qi-required-weight}
 B_{\mathrm{req}}=4+\log\mathcal K_\theta+\kappa,
 \qquad
 k_L=1+\left\lfloor\frac{(L-1)/2-1}{2}\right\rfloor.
\end{equation}
Connected-carrier compression and \eqref{qi-projected-size} give
\begin{equation}\label{qi-carrier-compression}
 |Y(H)|\le\chi2^d+(\chi2^d/k_L)|H|.
\end{equation}

\begin{proposition}[Conditional inherited interaction re-entry]
\label{qi-reentry-proposition}
Assume the exact source chart, \eqref{qi-inherited-norm},
\eqref{qi-domain-budget}, and the numerical inequality
\begin{equation}\label{qi-reentry-condition}
 \frac{\chi2^d}{k_L}B_{\mathrm{req}}\le a_0.
\end{equation}
Then, with
$V_{\max}=u N_{\mathrm{near}}L^d\epsilon$,
\begin{align}
 \sup_B\sum_{Y\ni B}e^{B_{\mathrm{req}}|Y|}\norm{Q_Y}
 &\le u e^{B_{\mathrm{req}}\chi2^d}
                       N_{\mathrm{near}}L^d\epsilon,              \label{qi-exponent-bound}\\
 \epsilon_{\mathrm{int}}
 &\le e^{V_{\max}}u e^{B_{\mathrm{req}}\chi2^d}
                       N_{\mathrm{near}}L^d\epsilon.              \label{qi-multiplier-bound}
\end{align}
In particular, for $V_{\max}\le\log2$ the first prefactor
in \eqref{qi-multiplier-bound} is at most two.
\end{proposition}
\begin{proof}
Equations \eqref{qi-increment-bound} and
\eqref{qi-projected-anchor} first give
\[
 \sup_B\sum_{Y\ni B}\norm{Q_Y}
       \le u N_{\mathrm{near}}L^d\epsilon=V_{\max}.
\]
Every individual exponent is therefore bounded by $V_{\max}$,
and the multiplier estimate is exactly
$v_Y\le e^{\norm{Q_Y}}\norm{Q_Y}\le e^{V_{\max}}\norm{Q_Y}$.
One cannot retain the smaller individual-label bound $u e_H$
after grouping many $H$ into one $Y$.

For the weighted bound, \eqref{qi-carrier-compression} and
\eqref{qi-reentry-condition} imply
\[
 e^{B_{\mathrm{req}}|Y(H)|}
       \le e^{B_{\mathrm{req}}\chi2^d}e^{a_0|H|}.
\]
Use \eqref{qi-projected-anchor} and anchor at the at most
$N_{\mathrm{near}}L^d$ current sites in neighboring blocks.
The input norm \eqref{qi-inherited-norm} proves
\eqref{qi-exponent-bound}. Connectedness gives
$\ell(Y)=|Y|-1$, so the target weight in
\eqref{qi-integration-norm} is at most
$e^{B_{\mathrm{req}}|Y|}$. The multiplier estimate proves
\eqref{qi-multiplier-bound}.
\end{proof}
Thus the edge-geometry weight is included, not only the
$e^{4|Y|}$ term. If another estimate needs a factor $|H|^q$,
require a strict margin
$\rho=a_0-(\chi2^d/k_L)B_{\mathrm{req}}>0$ and use
$\sup_{n\ge1}n^q e^{-\rho n}$. No extra reserve is assumed
without such a budget.

\subsection{Comparison of the connected logarithm coefficients}

Assume the remaining covariance and precision hypotheses of
Notebook18 hold on the independent background ball. Write its
common smallness threshold as $\eta_*$ and its canonical strict
connected coefficient estimate as
$\norm{\mathcal U}_{\mathrm{out}}\le C_{\mathrm{out}}\eta$
for $\eta\le\eta_*$. Here the output norm is the theorem's
connected support norm in the original independent background
variables, not merely the absolute scalar logarithm.

\begin{corollary}[Inherited contribution to the logarithm]
\label{qi-log-comparison}
Let $\eta_{\mathrm{base}}$ include the Gaussian tail and the
classical, measure and gauge interaction norms. If
\[
 \eta_{\mathrm{base}}\le\eta_*/4,
 \qquad\epsilon_{\mathrm{int}}\le\eta_*/4,
\]
then the canonical coefficients with and without the inherited
multipliers satisfy
\begin{equation}\label{qi-log-bound}
 \norm{\mathcal U_{\mathrm{full}}-
             \mathcal U_{\mathrm{base}}}_{\mathrm{out}}
       \le4C_{\mathrm{out}}\epsilon_{\mathrm{int}}.
\end{equation}
\end{corollary}
\begin{proof}
If $\epsilon_{\mathrm{int}}=0$ the statement is immediate.
Otherwise multiply only each inherited bounded multiplier
$\chi_Y(e^{Q_Y}-1)$ by a complex parameter $\lambda$, keeping
the base family fixed. At $\lambda=0$ and $1$ this gives the
two exact integrals. It does not replace $Q_Y$ by $\lambda Q_Y$.
Base and inherited factors with the same carrier are retained
as distinct labelled hyperedges, as permitted by the overlap
theorem; no cancellation or division by a base factor is used.
For $|\lambda|\le R_\lambda=\eta_*/(2\epsilon_{\mathrm{int}})
\ge2$, the total budget is at most $3\eta_*/4$.
The canonical coefficient vector $\mathcal U(\lambda)$ is
holomorphic in the output Banach space and bounded by
$C_{\mathrm{out}}\eta_*$. Summing its Cauchy coefficient
bounds gives
\[
 \norm{\mathcal U(1)-\mathcal U(0)}_{\mathrm{out}}
 \le\frac{C_{\mathrm{out}}\eta_*}{R_\lambda-1}
 \le4C_{\mathrm{out}}\epsilon_{\mathrm{int}}.
\]
\end{proof}
The fluctuation-independent sum $\sum_HE_H(a)$ remains outside
this comparison. It and the new connected coefficients are
evaluated at $a_*(w)$ afterwards and transported on a common
smaller domain by the transport theorem of Notebook 22.

\subsection{The open parameter and domain tests}

The tunable density estimate supplies
\begin{equation}\label{qi-dimension-tuning}
 \log\mathcal K_\theta\le\tfrac52 D_G\theta,
 \qquad \theta=\frac1{4D_G} \Longrightarrow\
 B_{\mathrm{req}}\le\frac{37}{8}+\kappa.
\end{equation}
Without this tuning, a cost $\log K_G=O(D_G)$ in the
sufficient interaction norm would contribute $O(L^{d-1})$
to \eqref{qi-reentry-condition}. Tuning removes that
particular artificial dimension loss. It does not remove
the smaller covariance-edge radius, the required $O(\theta)$
precision neighborhood, or the $2^{D_G/2}$ Gaussian tail
prefactor. The actual independent background ball must meet
those strengthened precision conditions. For the actual
linkwise color-ball cutoff
$\chi_B=\prod_{b\in B}\mathbf1_{\{|\xi_b|\le R\}}$,
the bad event is contained in $\{\norm{\xi_B}_2>R\}$.
The valid Euclidean-block tail therefore uses threshold $R$.
The good event only implies
$\norm{\xi_B}_2\le\sqrt{d(L^d-1)}R$; this opposite-direction
bound cannot improve the bad-event threshold. The chart
uses link supremum and uses the original $R$ directly.
No Euclidean block window is silently identified with the
source linkwise window.

For each fixed $L$, a logarithmic cutoff $R(g)$ can grow while
$gR(g)\to0$, so the chart margin and a fixed Gaussian tail
threshold can both be met by sufficiently small $g$. This
does not prove the required all-scale coupling recursion.

More fundamentally, the preceding integration supplies an
actual reserve $a_0=a_0(L)$. Existing fixed-$L$ results give
positivity, but no quantitative large-$L$ lower bound proving
\begin{equation}\label{qi-open-L-test}
 k_L a_0(L)\ge
       \chi2^d[4+\log\mathcal K_\theta(L)+\kappa(L)].
\end{equation}
Positivity alone is compatible with $a_0(L)=e^{-L}$ and
cannot establish this test. This is a missing estimate for
the actual available constants, not a counterexample to
Yang--Mills theory. A quantitative locality estimate as $L$
varies, or an integration theorem with a smaller fixed
carrier cost, is a concrete next target.

The domain is also not automatic. One needs
$\norm{a_*(w)}_\infty\le C_a(L)r\le b_-$, and coefficient
transport may require $C_a(L)r/b$ of order
$a_0/D_{\mathrm{bg}}$. No $L^{-1}$ contraction in link-angle
supremum has been proved; a pure-gauge stationary branch can
have crossing-link angle of order $w$. Repeated use of the
one-step radius losses is therefore not a stable domain
construction. The constants in \eqref{qi-multiplier-bound}
need not be contractive, either.

Finally, strict locality must be applied in the correct order.
On independent $a$, \eqref{qi-grouped-exponent} has the strict
joint $(a,\xi)$ carrier $Y$. Apply the Gaussian theorem on that
product domain, then substitute the actual $a_*(w)$ into its
holomorphic identity. Substitution first generally makes the
background dependence nonlocal in $w$; its rooted connecting
trees need their own support budget, and $Y$ alone is no
longer a strict $w$-support. A term independent of $w$ but
dependent on $\xi$ is still an interaction, not a vacuum
constant. Neither order removes the later radius and weight
losses.

The proposition and corollary establish a usable conditional
inherited-interaction module. The re-entry test, stable domains,
relevant and marginal extraction, running couplings, and
large-field and observable sectors remain separate obligations.

In the extended construction, Notebook 26 supplies the support-weight
test using independent fixed auxiliary cells. It re-expands the
classical precision and Gaussian normalization on those cells
and recovers a common positive quantum carrier exponent.
That argument resolves this particular norm-weight interface;
the domain and amplitude recurrence remains open.


\begin{thebibliography}{9}
\bibitem{prior24}
\emph{Toward a four-dimensional Yang--Mills construction},
research Notebooks 18--23 in this workspace.
The tunable strict-support Gaussian expansion, exact pivot
geometry, connected-carrier projection estimate, and
stationary activity transport are the principal inputs.
\bibitem{balaban11624}
T.~Ba\l aban,
\emph{Renormalization group approach to lattice gauge field
theories. II. Cluster expansions},
Commun. Math. Phys. \textbf{116} (1988), 1--22,
\href{https://scispace.com/pdf/renormalization-group-approach-to-lattice-gauge-field-4penm4k7hc.pdf}{original article copy}.
\bibitem{source24}
E.~Sussman and collaborators,
\href{https://github.com/ethanwsussman/Yang--Mills-LLM-experiment}{Yang--Mills-LLM-experiment},
deposited raw accounts at commit
\texttt{e392ef45013a581664bab76b90867170b30fe10e}.
\end{thebibliography}
\end{document}
